
\input texsis
\paper

\overfullrule=0pt
\def\frac#1#2{#1\over #2}
%
\def\bmatrix#1{\left[ \matrix{#1} \right]}
\def\be{{\beta}}
\def\toone{{t+1}}

%\def\frac#1#2{#1\over #2}
%\magnification=1200
\overfullrule=0pt
%


\centerline{\bf Homework 4}


\medskip
\noindent{\bf Notation alert:}
Please forgive us for recycling some notation. In particular, the state vector $x_t$ and control vector $u_t$ are recycled, being redefined as we move
from one of the problems below to the next problem.


\medskip
\noindent{\it Exercise 1} \quad {\bf Demand for labor}
\medskip
\noindent
A representative firm chooses a labor input process $\{n_{t+1}\}_{t=0}^\infty$ to maximize $E_0 \sum_{t=0}^\infty \beta^t \pi_t$,
where the discount factor $\beta \in (0,1)$, $E_0(\cdot)$ is a mathematical expectation operator conditional on date $0$ information, $n_0$
is a given initial condition,
$$ \pi_t = s_t n_t -f n_t^2 - d (n_{t+1} - n_t)^2 - w_t n_t, $$
  $s_t$ is an exogenous productivity shock, $w_t$ is an exogenous wage rate, $f > 0$, and $d > 0$.  The $\{s_t, w_t\}_{t=0}^\infty$ process is governed by
$$\eqalign{ z_{t+1}  & = A_{22} z_t + C_2 \epsilon_{t+1} \cr
            \bmatrix{s_t \cr w_t} & = S z_t } \leqno(1) $$
where $A_{22}$ is a matrix whose eigenvalues are bounded in modulus by $\beta^{-{\frac{1}{2}}}$ and $\{\epsilon_{t+1}\}_{t=0}^\infty$ is an i.i.d.~process
of ${\cal N}(0,I)$ random vectors; $z_0$ is a given initial condition.

\medskip
\noindent{\bf a.} Assume that the $\{s_t, w_t\}$ processes take the forms
$$\eqalign{ s_{t+1} & = \alpha_s + \rho_1 s_t + \sigma_s \epsilon_{1,t+1} \cr
           w_{t+1} & = \alpha_w + \gamma_1 w_t + \sigma_w \epsilon_{2,t+1} ,}$$
where $\{\epsilon_{i,t}\}$ are each i.i.d.~sequences of ${\cal N}(0,1)$ random variables.  Show how to map these processes into system
(1).

\medskip
\noindent{\bf b.} Please formulate the firm's optimum problem as a stochastic discounted optimal linear regulator problem. In particular, define the
firm's state vector $x_t$ and define the firm's control
  as $u_t = n_{t+1} - n_t$. Tell how to define the matrices $A_f, B_f, C_f$ that describe the firm's state-transition equations.
   Describe how to create the matrices $R_f, Q_f, N_f$ that define the firm's return function
$r_f (x_t, u_t) = x_t ' R_f x_t + u_t ' Q_f u_t + 2 u_t' N_f x_t$. Please describe the form of the firm's demand curve for labor.






\medskip
\noindent{\it Exercise 2.} \quad {\bf Labor supply}
\medskip
\noindent
A representative worker chooses a labor input process $\{n_{t+1}\}_{t=0}^\infty$ to maximize $E_0 \sum_{t=0}^\infty \beta^t (c_t + \upsilon_t)$
where  $c_t = w_t n_t$ is the worker's consumption rate, $\upsilon_t$ is the worker's utility from working,  $w_t$ is an exogenous wage rate, the discount factor $\beta \in (0,1)$, $E_0(\cdot)$ is a mathematical expectation operator conditional on date $0$ information, $n_0$
is a given initial condition,
$$ \upsilon_t =  - b_t n_t - h n_t^2 - q (n_{t+1} + n_t)^2 , $$
 $w_t$ is an exogenous wage, $h > 0$,  and $q > 0$.  The $\{b_t, w_t\}_{t=0}^\infty$ process is governed by
$$\eqalign{ z_{t+1}  & = A_{22} z_t + C_2 \epsilon_{t+1} \cr
            \bmatrix{b_t \cr w_t} & = S z_t } \leqno(1) $$
where $A_{22}$ is a matrix whose eigenvalues are bounded in modulus by $\beta^{-{\frac{1}{2}}}$ and $\{\epsilon_{t+1}\}_{t=0}^\infty$ is an i.i.d.~process
of ${\cal N}(0,I)$ random vectors; $z_0$ is a vector of given initial conditions; the initial labor supply $n_0$ is given.

\medskip

\noindent{\bf a.} Assume that the $\{b_t, w_t\}$ processes take the forms
$$\eqalign{ b_{t+1} & = \alpha_b + \rho_1 b_t + \sigma_b \epsilon_{1,t+1} \cr
           w_{t+1} & = \alpha_w + \gamma_1 w_t + \sigma_w \epsilon_{2,t+1} ,}$$
where $\{\epsilon_{i,t}\}$ are each i.i.d.~sequences of ${\cal N}(0,1)$ random variables.  Show how to map these processes into system
(1).
\medskip

\noindent{\bf b.} Please formulate the worker's optimum labor supply problem as a stochastic discounted optimal linear regulator problem.
In particular, define the
worker's state vector $x_t$ and define the worker's control
  as $u_t = n_{t+1} - n_t$. Tell how to define the matrices $A_w, B_w, C_w$ that describe the worker's state-transition equations.
   Describe how to create the matrices $R_w, Q_w, N_w$ that define the worker's return function
$r_w (x_t, u_t) = x_t ' R_w x_t + u_t ' Q_w u_t + 2 u_t' N_w x_t$.
    Please describe the form of the worker's supply curve for labor and represent it in state-space form.





\medskip
\noindent{\it Exercise 3} \quad {\bf A planning problem}
\medskip

\noindent A planner chooses a labor input process $\{n_{t+1}\}_{t=0}^\infty$ to maximize $E_0 \sum_{t=0}^\infty \beta^t (c_t + \upsilon_t)$, where
$$\eqalign{ c_t &  =  s_t n_t - f n_t^2 - d (n_{t+1} - n_t)^2 \cr
            \upsilon_t & =   - b_t n_t - h n_t^2 - q (n_{t+1} + n_t)^2 } $$
where $f > 0, h >0, d > 0, q > 0$, $s_t$ is a productivity shock,  $b_t$ is a preference shock, $c_t$ is a worker's consumption, and $\upsilon_t$
  is a worker's utility from working.  The shock processes $s_t, b_t$ are
governed by
$$\eqalign{ z_{t+1}  & = A_{22} z_t + C_2 \epsilon_{t+1} \cr
            \bmatrix{s_t \cr b_t} & = S z_t }$$
where $A_{22}$ is a matrix whose eigenvalues are bounded in modulus by $\beta^{-{\frac{1}{2}}}$ and $\{\epsilon_{t+1}\}_{t=0}^\infty$ is an i.i.d.~process
of ${\cal N}(0,I)$ random vectors; $z_0$ is a vector of given initial conditions; the initial  labor supply  $n_0$ is also given.

\medskip
\noindent{\bf a.} Please formulate the planner's  optimum labor supply problem as a stochastic discounted optimal linear regulator problem.
 in particular, define the
planner's state vector $x_t$ and define the planner's  control
  as $u_t = n_{t+1} - n_t$. Tell how to define the matrices $A_p, B_p, C_p$ that describe the planner's state-transition equations.
   Describe how to create the matrices $R_p, Q_p, N_p$ that define the planner's return function
$r_p (x_t, u_t) = x_t ' R_p x_t + u_t ' Q_p u_t + 2 u_t' N_p x_t$. Please describe the form of the planner's rule for setting $n_{t+1}$  and represent it in state-space form.

\medskip

\noindent
{\bf Hints:} Set $u_t =  n_{t+1} - n_t$, $x_t = \bmatrix{z_t \cr n_t}, C_p = \bmatrix{ C_2 \cr 0}$. The optimal decision rule is $u_{t} = - F_p x_t$ and the optimal closed loop system
is $x_{t+1} = (A_p - B_pF_p) x_t + C \epsilon_{t+1}. $
\medskip

\noindent {\it Exercise 4.} {\bf A big $K$, little $k$ application}

\medskip
\noindent For convenience, please rewrite the controlled state-space system that emerged from solving the planner's problem in exercise 3
as $ \tilde x_{t+1} = (\tilde A - \tilde B \tilde F) \tilde x_t + \tilde C \epsilon_{t+1}$, where $\tilde x_t = \bmatrix{ z_t \cr \tilde n_t}$
and $\tilde A = A_p, \tilde B = B_p, \tilde F = F_p, \tilde C = C_p$. Define a matrix $\tilde R_f $ to satisfy
$ \tilde x_t' \tilde R_f \tilde x_t = s_t \tilde n_t - f \tilde n_t^2$ and set $\tilde Q_f = - d$. ({\bf Quentin: warning: I might have put a wrong sign on $\tilde Q$. I want to pick off the net marginal product including the adjustment costs! I hope this is the sum I have defined in the following.})  Then define a putative wage rate $ w_t =  - 2 P_{w,n} \tilde x_t$, where $P_{w,n}$ is the last row of the matrix $P_w$ that satisfies the discrete Sylvester equation
$$ P_w = (\tilde R_f + F_p' \tilde Q_f F_p) + \beta A_p' P_w (A_p - B_p F_p).   $$
Please note the subscripts carefully: the matrices $\tilde R_f, \tilde Q_f$ have just been defined, while  matrices $A_p, B_p, F_p$ are from the planner's problem.
This discrete Sylvester equation can  be solved by iterating to convergence on $ P_w^{(j+1)} = (\tilde R_f + F_p' \tilde Q_f F_p) + \beta A_p' P_w^{(j)} (A_p - B_p F_p).   $
 %Let $\tilde P$ be the matrix in the quadratic form for the value function and let
% $\tilde \mu_t = 2  \tilde P \tilde x_t$ be  scaled Lagrange multipliers, where $\tilde \mu_t = \bmatrix{ \tilde \mu_{zt} \cr \tilde \mu_{nt}}$.


Now consider the representative firm's problem.  Suppose  that the wage rate $w_t$ and the shocks $s_t, b_t$ obey  the process
$$\eqalign{ \tilde x_{t+1} & = (\tilde A - \tilde B \tilde F) \tilde x_t + \tilde C \epsilon_{t+1} \cr
               \bmatrix{ s_t \cr b_t \cr w_t } & = \bmatrix{ S_s \cr S_b \cr - 2  P_{w,n} } \tilde x_t, } $$
where $\tilde P_{w,n}$ is the row of $ P_w$ corresponding to $\tilde n$ (i.e., the last row).
Formulate the representative firm's problem as a discounted dynamic programming problem in which the control is $\hat u_t = \hat n_{t+1} - \hat n_t$ and the state is
$\hat x_t = \bmatrix{ \tilde x_t \cr \hat n_t }$, where we use $(\hat \cdot)$ to denote variables chosen by the representative firm.
Find the firm's optimal decision rule in the form $\hat u_{t} = - \hat F \hat x_t$.  Then verify by computation
that if you set $\hat n_t = \tilde n_t$ it follows that
$\hat u_{t} = - \hat F \bmatrix{ \tilde z_t \cr \tilde n_t \cr \hat n_t } =  \tilde u_t = - \tilde F \bmatrix{ \tilde z_t \cr \tilde n_t} $.

\medskip

\noindent {\it Exercise 5.} {\bf Another big $K$, little $k$ application}

\medskip
\noindent {\bf To be repaired later.} Now consider the representative worker's labor supply problem. We'll recycle some notation, hopefully without causing confusion.  Again,
for convenience, please rewrite the controlled state-space system that emerged from  the planner's problem in exercise 3
as $ \tilde x_{t+1} = (\tilde A - \tilde B \tilde F) \tilde x_t + \tilde C \epsilon_{t+1}$ where $\tilde x_t = \bmatrix{ z_t \cr \tilde n_t}$.
 Let $\tilde P$ be the matrix in the quadratic form for the value function and let
 $\tilde \mu_t = 2 \tilde P \tilde x_t$ be the scaled Lagrange multipliers where $\tilde \mu_t = \bmatrix{ \tilde \mu_{zt} \cr \tilde \mu_{nt}}$.
 Suppose  that the wage rate $w_t$ and the shocks $s_t, b_t$ that the worker regards as exogenous obey  the process
$$\eqalign{ \tilde x_{t+1} & = (\tilde A - \tilde B \tilde F) \tilde x_t + \tilde C \epsilon_{t+1} \cr
               \bmatrix{ s_t \cr b_t \cr w_t } & = \bmatrix{ S_s \cr S_b \cr 2 \tilde P_n } \tilde x_t } $$
where $\tilde P_n$ is again the row of $\tilde P$ corresponding to $\tilde n$.
Formulate the representative worker's problem as a discounted dynamic programming problem in which the state is
$\hat x_t = \bmatrix{ \tilde x_t \cr \hat n_t }$, where we use $(\hat \cdot )$ to denote variables chosen by the worker.
Find the worker's optimal decision rule in the form $\hat u_{t} = - \hat F \hat x_t$.  Then verify by computation
that if you set $\hat n_t = \tilde n_t$ it follows that
$\hat u_{t} = - \hat F \bmatrix{ \tilde z_t \cr \tilde n_t \cr \hat n_t } = \tilde u_t = - \tilde F \bmatrix{ \tilde z_t \cr \tilde n_t} $.




\bye
